% chapters/pengembangan.tex
\chapter{Proses Pengembangan Dashboard}

\section{Pertanyaan Kunci}

Dashboard \textit{``Rokok Nomor Satu, Gizi Lain Waktu''} dirancang untuk menjawab tiga pertanyaan kunci:

\begin{enumerate}
  \item \textbf{Seberapa besar proporsi pengeluaran rokok dibandingkan pengeluaran gizi} di 34 provinsi Indonesia (berdasarkan data SUSENAS 2024)?
  \item \textbf{Apakah ada asosiasi} antara tingginya rasio pengeluaran rokok terhadap gizi dengan indikator kesehatan seperti prevalensi merokok, stunting balita, dan asupan protein per kapita?
  \item \textbf{Provinsi mana yang menjadi prioritas perhatian}, jika mempertimbangkan secara bersamaan indikator pengeluaran rokok, kemiskinan, stunting, dan protein?
\end{enumerate}

Ketiga pertanyaan ini membentuk alur narasi dashboard dari pandangan makro nasional, menuju diagnosis per provinsi, hingga rekomendasi prioritas kebijakan.

\section{Jenis Dashboard}

Dashboard ini termasuk kategori \textbf{Dashboard Analitis} (\textit{Analytical Dashboard}). Alasannya:

\begin{itemize}
  \item Data disajikan bukan sekadar untuk memantau kondisi saat ini, melainkan untuk mengeksplorasi pola, asosiasi, dan hipotesis.
  \item Pengguna dapat melakukan \textit{drill-down} dari tampilan nasional ke diagnosa per provinsi.
  \item Tersedia simulasi skenario (\textit{what-if}) yang memungkinkan eksplorasi realokasi pengeluaran.
  \item Disediakan konteks multi-indikator (pengeluaran, kesehatan, demografi) untuk mendukung pengambilan keputusan berbasis data.
\end{itemize}

\section{Target Pengguna}

Dashboard ditujukan kepada tiga segmen pengguna utama:

\begin{enumerate}
  \item \textbf{Peneliti dan analis kebijakan kesehatan masyarakat} -- yang membutuhkan eksplorasi asosiasi antara kebiasaan konsumsi rumah tangga dengan indikator gizi dan kesehatan.
  \item \textbf{Mahasiswa dan akademisi} -- yang mempelajari visualisasi data, ekonomi kesehatan, atau epidemiologi sosial.
  \item \textbf{Jurnalis dan komunikator data} -- yang ingin memahami konteks regional pengeluaran rokok di Indonesia.
\end{enumerate}

Semua segmen pengguna ini diasumsikan memiliki literasi data dasar dan tidak memerlukan penjelasan teknis mendalam untuk memahami grafik yang disajikan.

\section{Pemilihan Jenis Grafik}

Pemilihan jenis grafik didasarkan pada karakteristik data dan tujuan komunikatif masing-masing bagian dashboard:

\begin{table}[H]
\centering
\caption{Pemilihan jenis grafik berdasarkan data dan tujuan}
\renewcommand{\arraystretch}{1.3}
\begin{tabularx}{\textwidth}{@{}lXX@{}}
\toprule
\textbf{Visualisasi} & \textbf{Data} & \textbf{Justifikasi} \\
\midrule
Choropleth Map & Rasio rokok/gizi per provinsi (quantitative-ratio) & Distribusi spasial 38 provinsi dalam satu pandangan; warna berurutan (putih ke merah) mencerminkan intensitas rasio \\
Donut Chart & Komposisi pengeluaran rokok vs.\ 5 komponen gizi & Proporsi sederhana antar kategori; efektif untuk impresi visual ``berapa besar irisan rokok'' \\
Horizontal Bar Chart & Ranking 10 provinsi tertinggi (rokok \% gizi) & Perbandingan nilai antar kategori; label langsung terbaca \\
Butterfly Chart & Data karakteristik (gender, pendidikan, pekerjaan, dll.) & Membandingkan dua variabel berlawanan (merokok kiri, gizi normal kanan) dalam satu sumbu \\
Scatter Plot (Kompas) & Rokok/gizi vs.\ protein per kapita, seluruh provinsi & Menampilkan posisi relatif provinsi di dua dimensi sekaligus; ukuran titik merepresentasikan populasi \\
Sankey Diagram & Simulasi aliran dana hemat rokok ke komponen gizi & Memvisualisasikan aliran kuantitatif dari satu sumber ke banyak tujuan \\
Leaderboard Bar & Struktur pengeluaran satu provinsi & Perbandingan 6 komponen (rokok + 5 gizi) secara langsung \\
\bottomrule
\end{tabularx}
\end{table}

\section{Persiapan Data}

\subsection{Sumber Data}

Data yang digunakan dalam dashboard ini berasal dari tiga sumber utama:

\begin{enumerate}
  \item \textbf{SUSENAS (Survei Sosial Ekonomi Nasional) 2024 -- BPS}: Data pengeluaran per kapita untuk 6 komoditas (rokok, sayur, ikan, telur/susu, daging, buah) di 34 provinsi Indonesia (tidak termasuk 4 provinsi pemekaran Papua 2022 yang belum tercakup).
  \item \textbf{SKI (Survei Kesehatan Indonesia) 2023 -- Kemenkes}: Data prevalensi merokok (umur 10--18 tahun dan 10+ tahun), stunting balita (0--59 bulan), pola konsumsi pangan (ikan, telur, susu), harga rokok, dan rata-rata batang rokok per hari di 38 provinsi.
  \item \textbf{Data status gizi berdasarkan status ekonomi -- SKI 2023}: Data gizi normal anak 5--12 tahun per kuintil ekonomi, digunakan untuk visualisasi butterfly chart dimensi status ekonomi.
\end{enumerate}

\subsection{Pembersihan dan Transformasi Data}

Proses pembersihan dan transformasi dilakukan melalui pipeline Python (\opus{scripts/prepare\_data.py}):

\begin{itemize}
  \item \textbf{Normalisasi nama provinsi}: Semua nama provinsi dari berbagai sumber diseragamkan menggunakan fungsi \opus{normalize\_province\_name()} yang menangani alias (DKI Jakarta $\to$ Jakarta Raya, Kepulauan Bangka Belitung, dll.) dan inkonsistensi kapitalisasi.
  \item \textbf{Pembuatan variabel turunan}: Dihitung rasio \opus{rokok\_pct\_of\_gizi} (rokok / total gizi $\times$ 100), \opus{rokok\_pct\_of\_sayur}, dan \opus{rokok\_pct\_of\_daging}.
  \item \textbf{Penggabungan dataset}: Data SUSENAS, SKI provinsi, dan SKI karakteristik digabungkan menggunakan kunci \opus{province} yang sudah dinormalisasi.
  \item \textbf{Penanganan data tidak lengkap}: 4 provinsi Papua baru (Papua Barat Daya, Papua Pegunungan, Papua Selatan, Papua Tengah) ditandai \opus{has\_expenditure\_data = False} dan ditampilkan dalam mode \textit{greyed-out} di peta.
  \item \textbf{Regresi linear}: Model regresi linear sederhana dilatih lintas 34 provinsi untuk hubungan \opus{gizi\_total} $\to$ \opus{stunting\_pct} dan \opus{gizi\_total} $\to$ \opus{protein\_per\_capita}, disimpan sebagai \opus{regression\_models.json} untuk digunakan di simulator kebijakan.
  \item \textbf{Data butterfly}: Untuk setiap dimensi karakteristik (gender, pendidikan, pekerjaan, tempat tinggal, status ekonomi), dibuat file CSV tersendiri yang memetakan label tampilan ke nilai prevalensi merokok dan gizi normal.
\end{itemize}

\subsection{Pengelompokan Data}

Data final disimpan dalam beberapa file di direktori \opus{data/clean/}:

\begin{itemize}
  \item \opus{master\_profile.csv} -- 38 provinsi dengan seluruh indikator tergabung
  \item \opus{butterfly\_\{dimensi\}.csv} -- 5 file untuk masing-masing dimensi karakteristik
  \item \opus{dim\_province.csv} -- metadata geografis (koordinat, geo feature ID, region)
  \item \opus{regression\_models.json} -- koefisien dan R\textsuperscript{2} model regresi
\end{itemize}

\section{Struktur Informasi}

Dashboard diorganisasikan dalam dua halaman utama dengan alur naratif yang progresif:

\textbf{Halaman 1 -- Indonesia (Pandangan Nasional)}

Pengguna disambut dengan gambaran besar: KPI strip nasional (rata-rata rokok/gizi, prevalensi perokok, stunting balita), peta choropleth 38 provinsi, donut komposisi pengeluaran nasional, ranking bar 10 provinsi tertinggi, dan butterfly chart karakteristik sosial-ekonomi. Dari sini, pengguna bisa mengklik provinsi di peta untuk masuk ke diagnosa per provinsi.

\textbf{Halaman 2 -- Provinsi (Diagnosa Lokal)}

Menampilkan profil lengkap satu provinsi: dua baris KPI (pengeluaran, kesehatan, ranking), leaderboard struktur pengeluaran, kompas kuadran posisi provinsi, dan simulator kebijakan interaktif (slider realokasi + Sankey diagram).

Dua kontrol global tersedia di atas dashboard:
\begin{itemize}
  \item \textbf{Dropdown Metrik}: memilih metrik peta (Rokok \% Gizi Total, Rokok \% Sayuran, Rokok \% Daging, Stunting Balita)
  \item \textbf{Dropdown Region}: menyaring tampilan berdasarkan pulau/wilayah (Seluruh Indonesia, Jawa, Sumatera, Kalimantan, Sulawesi, Bali-Nusa Tenggara, Maluku-Papua)
\end{itemize}

\section{Rancangan Tampilan Visual}

\subsection{Skema Warna}

Dashboard menggunakan dark theme berbasis coklat hangat (\opus{bg\_app: \#26211F}) dengan palet warna yang bermakna semantik, didefinisikan dalam satu sumber kebenaran (\opus{tokens.py}):

\begin{itemize}
  \item \textbf{Merah tembakau} (\opus{\#C0272D}): semua elemen terkait rokok dan tembakau
  \item \textbf{Hijau gizi} (\opus{\#3DD68C}): semua elemen terkait gizi dan protein
  \item \textbf{Emas} (\opus{\#DAA520}): highlight, rata-rata nasional, dan elemen navigasi aktif
  \item \textbf{Teal gizi} (\opus{\#2DD4BF}): garis tren gizi normal pada dual-axis chart
  \item \textbf{Maroon daging} (\opus{\#A03A2C}): komponen daging -- dibedakan dari merah tembakau
  \item \textbf{Abu-abu hangat} (\opus{\#99948F}): teks sekunder, label grafik, dan narasi
\end{itemize}

Skala choropleth menggunakan gradasi abu-abu ke merah tua ke merah cerah (\opus{\#4D4D4D} $\to$ \opus{\#5C0000} $\to$ \opus{\#C0272D} $\to$ \opus{\#F73227}), sehingga intensitas merah secara intuitif mencerminkan tingginya rasio rokok terhadap gizi. Untuk metrik protein/stunting, skala divergen merah $\to$ abu $\to$ hijau digunakan.

\subsection{Tipografi}

Tiga keluarga font digunakan secara konsisten, dimuat via Google Fonts:
\begin{itemize}
  \item \textbf{Oswald}: angka KPI besar, judul halaman, dan \textit{section kicker} -- display, tegas
  \item \textbf{Inter}: body text, label, narasi, dan kontrol UI -- modern, legibel
  \item \textbf{JetBrains Mono}: nilai monospace (kode, ID) -- konsistensi spasi
\end{itemize}

\subsection{Prinsip Layout}

\begin{itemize}
  \item \textbf{F-pattern}: elemen terpenting (KPI, peta) di kiri atas
  \item \textbf{Hierarki visual}: ukuran dan warna menciptakan urutan baca yang jelas
  \item \textbf{Whitespace}: \textit{card} dengan border tipis memisahkan area tanpa terasa penuh
  \item \textbf{Section kicker}: label kecil berwarna emas sebagai penanda seksi
\end{itemize}

\section{Desain Awal Dashboard}

Desain awal disusun dalam tiga iterasi:

\begin{enumerate}
  \item \textbf{Iterasi 1 -- Wireframe}: Layout kasar dua halaman (Indonesia dan Provinsi), menentukan posisi peta, KPI, dan bar chart. Diputuskan untuk menggunakan Plotly Dash karena fleksibilitas callback interaktif.
  \item \textbf{Iterasi 2 -- Prototipe Fungsional}: Implementasi visualisasi utama (choropleth, donut, bar, butterfly) dengan data nyata. Ditemukan masalah ketidakcocokan nama provinsi antara GeoJSON dan dataset -- diselesaikan dengan modul normalisasi.
  \item \textbf{Iterasi 3 -- Penyempurnaan Desain}: Penambahan dark theme, palet warna semantik, tooltip informatif, dan mekanisme \textit{fallback} peta (point map jika GeoJSON tidak lengkap).
\end{enumerate}

\section{Uji Coba Awal}

Uji coba dilakukan terhadap 5 responden yang merepresentasikan target pengguna dashboard. Responden diminta untuk menjawab tiga pertanyaan evaluasi:

\begin{enumerate}
  \item Apakah visualisasi mudah dipahami?
  \item Apakah \textit{insight} dapat diperoleh dari dashboard?
  \item Apakah ada elemen yang membingungkan?
\end{enumerate}

\subsection{Profil Responden}

\begin{table}[H]
\centering
\caption{Profil responden uji coba}
\renewcommand{\arraystretch}{1.3}
\begin{tabularx}{\textwidth}{@{}clXc@{}}
\toprule
\textbf{No.} & \textbf{Latar Belakang} & \textbf{Relevansi} & \textbf{Literasi Data} \\
\midrule
1 & Mahasiswa Teknik Informatika & Target akademisi & Tinggi \\
2 & Mahasiswa Teknik Informatika & Target akademisi & Tinggi \\
3 & Mahasiswa non-teknis & Pengguna umum & Sedang \\
4 & Peneliti kebijakan kesehatan & Target utama & Tinggi \\
5 & Mahasiswa non-teknis & Pengguna umum & Rendah \\
\bottomrule
\end{tabularx}
\end{table}

\subsection{Hasil Uji Coba}

Temuan utama dari uji coba:

\begin{itemize}
  \item \textbf{Peta choropleth}: Semua responden memahami makna gradasi warna. Responden 5 (literasi rendah) awalnya bingung dengan 4 provinsi greyed-out -- perlu penjelasan teks.
  \item \textbf{Butterfly chart}: 3 dari 5 responden memerlukan waktu lebih lama untuk memahami cara baca dua sisi. Label ``Merokok (kiri)'' dan ``Gizi Normal (kanan)'' perlu diperjelas.
  \item \textbf{Simulator kebijakan}: Semua responden menyukai fitur ini, namun 2 responden khawatir dashboard terkesan membuat klaim kausal. Perlu penambahan disclaimer yang lebih jelas.
  \item \textbf{Donut chart}: Semua responden langsung memahami komposisi tanpa penjelasan.
  \item \textbf{Dark theme}: Responden mahasiswa menilai positif; 1 responden peneliti menyarankan opsi \textit{light mode}.
\end{itemize}

\section{Penyempurnaan Setelah Uji Coba}

Berdasarkan hasil uji coba, beberapa penyempurnaan dilakukan:

\begin{enumerate}
  \item \textbf{Penambahan teks disclaimer} di bawah butterfly chart dan simulator: ``Estimasi berdasarkan regresi linear... Ini adalah asosiasi statistik, bukan prediksi kausal.''
  \item \textbf{Penambahan narasi otomatis} di kartu Insight halaman Indonesia, yang secara dinamis menjelaskan jumlah provinsi aktif, rata-rata metrik, dan cara interaksi dengan peta.
  \item \textbf{Perbaikan label butterfly chart}: Ditambahkan penjelasan ``Sisi kiri: \% perokok harian. Sisi kanan: \% gizi normal balita.''
  \item \textbf{Perbaikan tooltip peta}: Tooltip diperkaya dengan informasi protein, kemiskinan, dan stunting -- tidak hanya rokok dan gizi.
  \item \textbf{Penambahan badge risiko} pada halaman provinsi: ``Risiko Tinggi'' (merah), ``Risiko Sedang'' (kuning), atau ``Risiko Rendah'' (hijau) berdasarkan nilai \linebreak\opus{rokok\_pct\_of\_gizi}.
  \item \textbf{Mekanisme fallback peta}: Jika GeoJSON tidak lengkap, dashboard otomatis beralih ke bubble map berbasis koordinat agar tidak ada provinsi yang hilang.
\end{enumerate}

\section{Evaluasi Hasil Akhir}

\subsection{Fungsi}

Dashboard berfungsi penuh dengan semua fitur interaktif bekerja sebagaimana dirancang:
\begin{itemize}
  \item Klik provinsi di peta $\to$ navigasi ke halaman provinsi
  \item Dropdown metrik dan region $\to$ pembaruan peta, bar, dan KPI secara sinkron
  \item Slider realokasi $\to$ pembaruan estimasi stunting, protein, tabungan, dan Sankey secara real-time
  \item Dropdown butterfly $\to$ pembaruan chart karakteristik seketika
\end{itemize}

\subsection{Kemudahan Penggunaan}

\begin{itemize}
  \item Alur navigasi intuitif: dari pandangan nasional, klik untuk masuk ke detail provinsi
  \item Kontrol global (metrik + region) selalu terlihat di atas halaman
  \item Tombol ``Kembali ke Indonesia'' tersedia di halaman provinsi
  \item Tooltip informatif pada semua elemen grafik
\end{itemize}

\subsection{Estetika}

\begin{itemize}
  \item Dark theme coklat hangat (\opus{\#26211F}) konsisten di seluruh halaman, kartu (\opus{\#2D2B28}), dan header (\opus{\#1E1A18})
  \item Hierarki tipografi jelas: Oswald display (KPI besar) $\to$ Oswald heading $\to$ Inter body
  \item Palet warna semantik satu sumber (\opus{tokens.py}): merah tembakau = rokok, hijau = gizi, emas = highlight/navigasi, maroon = daging
  \item Warna teks memenuhi standar kontras WCAG AA: teks muted \opus{\#8C8984} memberikan rasio $\geq$4.5:1 pada latar gelap
  \item Whitespace memadai: padding kartu 16--20px, \textit{gap} antar kolom 16px
\end{itemize}
